
\documentclass[11pt,a4paper]{article}
\usepackage[utf8]{inputenc}
\usepackage[T1]{fontenc}
\usepackage[spanish]{babel}
\usepackage{geometry}
\usepackage{graphicx}
\usepackage{hyperref}
\usepackage{xcolor}
\usepackage{tabularx}
\geometry{margin=2cm}
\setlength{\parindent}{0pt}
\setlength{\parskip}{0.55em}
\sloppy
\definecolor{tcTeal}{HTML}{0E7C74}
\definecolor{tcDark}{HTML}{101817}
\definecolor{tcMuted}{HTML}{526A65}
\definecolor{tcWarnBg}{HTML}{FFF7E5}
\definecolor{tcWarnBorder}{HTML}{C99123}
\hypersetup{colorlinks=true, urlcolor=tcTeal, linkcolor=tcTeal}
\begin{document}

\begin{center}
{\huge\bfseries Informe AI Analyst: XAGEUR.r|XAGUSD.r H1 correlation}\\[0.25em]
{\large Reporte read-only para revision humana}
\end{center}

\vspace{0.2em}
\hrule height 0.7pt
\vspace{0.9em}

\begin{tabularx}{\textwidth}{>{\bfseries}p{3.2cm}X}
Paquete & correlation H1 XAGEUR.r XAGUSD.r d67ee4c9 \\
Setup & XAGEUR.r|XAGUSD.r H1 correlation \\
Modelo & gpt-5.5 \\
Thinking & medium \\
Macro & not requested \\
Estado & reviewed \\
\end{tabularx}

\vspace{0.75em}
\fcolorbox{tcWarnBorder}{tcWarnBg}{%
\begin{minipage}{0.94\textwidth}
\textbf{Limite de uso.} Este informe no es una senal, no autoriza MT5 y no ejecuta ordenes. Sirve para priorizar revision humana y documentar el contexto.
\end{minipage}}

\vspace{0.4em}

\section*{Resumen}
Revision read-only de correlacion H1 entre XAGEUR.r y XAGUSD.r. La relacion estadistica es extremadamente alta y positiva: Pearson 0.996903, Spearman 0.992718, Kendall 0.937668 y distancia de correlacion 0.994052 sobre 1999 observaciones. El grafico respalda visualmente una nube casi lineal, con pocos puntos alejados pero sin ruptura evidente del vinculo principal.

\section*{Lectura tecnica}
El paquete mide retornos logaritmicos cierre a cierre en H1 para plata cotizada frente a EUR y plata cotizada frente a USD, con muestra entre 2025-11-11 22:00:00 y 2026-03-17 09:00:00. La lectura principal es de co-movimiento fuerte, no de oportunidad operativa. En el radar de mercado mas reciente, ambos instrumentos aparecen dentro de metales, con tendencia bajista alineada en H1, H4 y D1, estado de sobreventa en H1/H4 y volatilidad H1 superior a su mediana historica.



\section*{Grafico de referencia}
\begin{center}
\includegraphics[width=0.96\textwidth]{\detokenize{ai_analyst_report_chart.png}}
\end{center}


\section*{Lectura de confluencias}
Todas las metricas de dependencia coinciden en una relacion positiva muy alta, no solo Pearson; esto reduce la posibilidad de que la lectura dependa de una unica metrica lineal; La correlacion movil de 120 periodos sigue elevada: Pearson 0.994939, Spearman 0.991958, Kendall 0.928291 y distancia de correlacion 0.994052, con variacion positiva frente al periodo previo; La dispersion visual muestra retornos de ambos activos moviendose casi sobre la misma diagonal, coherente con que ambos representan plata con distinta divisa de cotizacion; y El contexto de mercado para XAGEUR.r y XAGUSD.r es similar: ambos figuran como metales con alineacion bajista H1-H4-D1 y lectura de sobreventa en H1/H4.

\section*{Contradicciones y cautelas}
La muestra global de correlacion termina el 2026-03-17 09:00:00, mientras el radar de mercado para XAGEUR.r y XAGUSD.r usa datos de junio de 2026; por tanto, correlacion historica y contexto actual no pertenecen exactamente a la misma ventana temporal; XAGEUR.r tiene hora de radar 2026-06-05 15:00:00 y XAGUSD.r 2026-06-05 22:15:00, lo que introduce una diferencia temporal en la lectura comparativa reciente; La correlacion elevada no implica que el spread entre ambos sea estable ni operable: falta informacion directa de EURUSD, diferenciales de cotizacion, coste, liquidez y estabilidad fuera de muestra; y La distancia de correlacion movil declara ventana 120 pero obs 80, por lo que conviene revisar si hay datos faltantes o una limitacion metodologica en ese calculo.

\section*{Riesgos principales}
Riesgo de interpretar una correlacion casi perfecta como senal direccional. El paquete esta marcado como estudio, sin ejecucion ni generacion de senales; Riesgo de falsa robustez por activos economicamente solapados: ambos dependen de la plata, y la diferencia de divisa puede quedar subrepresentada si no se analiza explicitamente EURUSD; Riesgo de desfase temporal: la correlacion principal acaba en marzo de 2026, mientras el radar usado como contexto esta generado en junio de 2026; Riesgo de concentracion: en un evento fuerte de metales, ambos activos pueden moverse juntos y no aportar diversificacion efectiva; y Riesgo de outliers: la nube incluye retornos extremos, por ejemplo caidas cercanas a -3.4\% y movimientos positivos superiores a 3\%, que deberian auditarse por impacto en las metricas.

\section*{Contexto macro y noticias}
No se solicito lectura macro para este paquete. El informe se limita al contexto tecnico y a los datos estructurados del Trading Center.

\section*{Conclusion operativa y financiera}
Conclusion operativa y financiera: el analisis de correlacion queda como lectura de relacion entre activos, con prioridad de revision AI 2/5 y riesgo macro no solicitado. La conclusion debe usarse para detectar dependencia, diversificacion aparente o concentracion de riesgo entre instrumentos, no para validar un setup por si sola. Si la relacion es inestable, cambia por timeframe o se apoya en una ventana corta, debe tratarse como contexto auxiliar antes de revisar exposicion o pares comparables. En todos los casos, el informe documenta seguimiento y revision humana; no autoriza ejecucion, MT5 ni envio de ordenes.

\section*{Comprobaciones humanas}
\begin{itemize}
  \item Verificar con la serie completa si la correlacion se mantiene en una ventana actualizada hasta junio de 2026.
  \item Revisar rolling correlations por subperiodos y eventos para detectar rupturas temporales del vinculo.
  \item Comparar la relacion XAGEUR.r-XAGUSD.r contra EURUSD para separar efecto plata y efecto divisa.
  \item Auditar los puntos extremos de la muestra y confirmar que no son gaps, errores de datos o diferencias de horario entre simbolos.
  \item Comprobar costes, liquidez, spread y sincronizacion de velas antes de usar esta relacion en cualquier estudio posterior no operativo.
\end{itemize}

\section*{Fuentes}
\textbf{Fuentes locales}\par
\begin{tabularx}{\textwidth}{p{4.2cm}X}
setup\_context.metrics & Fuente local del paquete de revision \\
market\_context.correlation\_pair & Fuente local del paquete de revision \\
chart\_layers\_tail & Fuente local del paquete de revision \\
ohlc\_window\_tail & Fuente local del paquete de revision \\
source\_manifest.sources & Fuente local del paquete de revision \\
chart.png & Imagen renderizada del grafico \\
\end{tabularx}

\end{document}
